% from ega4-14-fr.tex
\oldpage[IV-3]{217}

On peut se borner au cas où $X=U$ est un voisinage ouvert affine de $x$ tel que
toutes les composantes irréductibles de $f^{-1}(y)$ contiennent $x$. Ces
composantes correspondent aux idéaux premiers minimaux de
$\mathcal O_x/\mathfrak m_y\mathcal O_x$, et par hypothèse ces idéaux sont
distincts de $\mathfrak m_x/\mathfrak m_y\mathcal O_x$ (Bourbaki,
\emph{Alg. comm.}, chap.~II, \S~1, n\textsuperscript{o}~1, prop.~2)~; pour obtenir
un $g\in\Gamma(U,\mathcal O_X)$ vérifiant les conditions de l'énoncé, il suffit de
prendre $g\in\mathfrak j_x$ tel que l'image de $g$ dans $\mathfrak m_x$
n'appartienne à aucun des idéaux premiers précédents. D'ailleurs, on a
$\operatorname{codim}(X',X)\leq 1$ (5.1.8), et comme $X'$ ne contient aucune des composantes
irréductibles de $f^{-1}(y)$ et que celles-ci sont de dimension $e$, on a
($\mathbf 0$, 14.2.2.2) $\dim_x(X'\cap f^{-1}(y))\leq e-1$. Il suffit alors
d'appliquer (14.5.1).

\begin{proposition}[14.5.3]
\label{IV.14.5.3-fr}
Soient $Y$ un préschéma irréductible localement noethérien, $f:X\to Y$ un
morphisme localement de type fini, $x$ un point de $X$. On suppose que $f$ est
équidimensionnel au point $x$ et que $x$ est fermé dans $f^{-1}(f(x))$. Alors il
existe une partie irréductible $X'$ de $X$, localement fermée dans $X$, contenant
$x$ et telle que la restriction $X'\to Y$ de $f$ (où $X'$ est le sous-préschéma
réduit de $X$ ayant $X'$ pour espace sous-jacent) soit un morphisme quasi-fini et
dominant.
\end{proposition}

En effet, avec les notations de (14.5.2), l'hypothèse que $x$ est fermé dans
$f^{-1}(y)$ entraîne que $x$ n'est pas point maximal de $X'\cap f^{-1}(y)$ tant
que $e-1\geq 1$. Il suffit donc d'appliquer (14.5.2) en raisonnant par récurrence
descendante sur $e=\dim_x(f^{-1}(f(x)))$ jusqu'à ce que l'on arrive à $e=1$~;
l'application de (14.5.2) dans ce dernier cas donne un $X'$ tel que
$X'\cap f^{-1}(f(x))$ soit noethérien et de dimension 0, donc fini et discret~;
comme alors $X'$ est équidimensionnel sur $Y$, $X'\cap f^{-1}(f(x'))$ est de
dimension 0 pour tout $x'\in X'$, ce qui entraîne que la restriction $X'\to Y$ de
$f$ est un morphisme quasi-fini (\textbf{II}, 6.2.2).

\begin{corollary}[14.5.4]
\label{IV.14.5.4-fr}
Sous les hypothèses de (14.5.3), supposons en outre que $Y=\operatorname{Spec}(A)$, où $A$ est
un anneau local noethérien intègre et complet, et que $y=f(x)$ soit l'unique point
fermé de $Y$. Alors, il existe un anneau local intègre $A'$, contenant $A$, qui est
une $A$-algèbre finie et possède la propriété suivante~: si l'on pose
$Y'=\operatorname{Spec}(A')$ et $X'=X\times_Y Y'$, il existe une $Y'$-section $h:Y'\to X'$ telle
que le morphisme composé $Y'\xrightarrow{h}X'\to X$ soit une immersion dont
l'image contient $x$.
\end{corollary}

Remplaçant au besoin $X$ par un sous-préschéma réduit irréductible de $X$, on peut,
en vertu de (14.5.3), se borner au cas où le morphisme $f$ est déjà quasi-fini et
dominant. Utilisant (\textbf{II}, 6.2.5), on en déduit que $\mathcal O_x=A'$ est un
anneau intègre qui est une $A$-algèbre finie, et que $X$ est somme disjointe du
sous-préschéma fermé $Y'=\operatorname{Spec}(A')$ et d'un sous-préschéma $Y''$~; le schéma $Y'$
répond à la question, le morphisme composé $Y'\to X'\to X$ n'étant autre que le
morphisme canonique $\operatorname{Spec}(\mathcal O_x)\to X$.

\begin{remark}[14.5.5]
\label{IV.14.5.5-fr}
Si on n'exige pas que dans l'énoncé de (14.5.4), le morphisme $Y'\to X$ soit une
immersion, on peut supposer que $A'$ est en outre \emph{intégralement clos}~: il
suffit en effet de remplacer $A'$ par sa \emph{clôture intégrale} $A_1$, car on
sait ($\mathbf 0$, 23.1.5) que $A_1$ est un $A$-module de type fini.
\end{remark}

\begin{proposition}[14.5.6]
\label{IV.14.5.6-fr}
Soient $Y$ un préschéma localement noethérien, irréductible, régulier et de
dimension 1, $f:X\to Y$ un morphisme localement de type fini, $y$ un point de $Y$.
Les conditions suivantes sont équivalentes~:
\begin{enumerate}[label=\emph{\alph*)}]
  \item $f_{\mathrm{red}}$ est plat en tout point de $f^{-1}(y)$.

  